\ifdefined\gkver\else\def\gkver{student}\fi
\documentclass[\gkver]{gaokaozhenti}
\begin{document}

\chapter{2011年北京理综}
%% paperType: 理综物理部分

\begin{enumerate}
\item
%% number: 13
%% paperName: 2011年北京理综第 13 题 6 分
%% typeId: 1
%% type: 单选题
%% chapter: 原子和原子核
%% point: 天然放射性
%% method: 无
%% score: 6
%% degree: 650
%% duplicateId: 0
%% body:
表示放射性元素碘 131（\ce{^{131}_{53}I}）$\beta$ 衰变的方程是\xzanswer{B}

\fourchoices[answer=B]
{\ce{^{131}_{53}I -> ^{127}_{51}Sb + ^{4}_{2}He}}
{\ce{^{131}_{53}I -> ^{131}_{54}Xe + ^{0}_{-1}e}}
{\ce{^{131}_{53}I -> ^{130}_{53}I + ^{1}_{0}n}}
{\ce{^{131}_{53}I -> ^{130}_{52}Te + ^{1}_{1}H}}
%% answer: B
%% memo:
\memoanswer{
由核反应过程电荷数守恒和质量数守恒及 $\beta$ 衰变释放可判断 B 正确。
}

\item
%% number: 14
%% paperName: 2011年北京理综第 14 题 6 分
%% typeId: 1
%% type: 单选题
%% chapter: 光学
%% point: 光的干涉
%% method: 无
%% score: 6
%% degree: 650
%% duplicateId: 0
%% body:
如图所示的双缝干涉实验，用绿光照射单缝 S 时，在光屏 P 上观察到干涉条纹。要得到相邻条纹间距更大的干涉图样，可以\xzanswer{C}

\onepicture[w=6cm]{figs/14.png}

\fourchoices[answer=C]
{增大 $S_{1}$ 与 $S_{2}$ 的间距}
{减小双缝屏到光屏的距离}
{将绿光换为红光}
{将绿光换为紫光}
%% answer: C
%% memo:
\memoanswer{
相邻条纹的间距 $\Delta x=\frac{D}{d}\lambda$，其中 D 指双缝屏到光屏的距离，d 指双缝之间的距离，$\lambda$ 指单色光的波长，故 A 错误；减小双缝屏到光屏的距离即减小 D，则 $\Delta x$ 减小，故 B 错误；将绿光换为红光，波长变长，相邻条纹间距变大，故 C 正确；同理可知 D 错误。
}

\item
%% number: 15
%% paperName: 2011年北京理综第 15 题 6 分
%% typeId: 1
%% type: 单选题
%% chapter: 曲线运动
%% point: 人造地球卫星
%% method: 无
%% score: 6
%% degree: 650
%% duplicateId: 0
%% body:
由于通讯和广播等方面的需要，许多国家发射了地球同步轨道卫星，这些卫星的\xzanswer{A}

\fourchoices[answer=A]
{质量可以不同}
{轨道半径可以不同}
{轨道平面可以不同}
{速率可以不同}
%% answer: A
%% memo:
\memoanswer{
发射地球同步卫星对卫星的质量没有要求，而轨道半径、轨道平面和速率必须相同，故 A 正确。
}

\item
%% number: 16
%% paperName: 2011年北京理综第 16 题 6 分
%% typeId: 1
%% type: 单选题
%% chapter: 机械波
%% point: 波长频率波速
%% method: 无
%% score: 6
%% degree: 650
%% duplicateId: 0
%% body:
介质中有一列简谐机械波传播，对于其中某个振动质点，\xzanswer{D}

\fourchoices[answer=D]
{它的振动速度等于波的传播速度}
{它的振动方向一定垂直于波的传播方向}
{它在一个周期内走过的路程等于一个波长}
{它的振动频率等于波源的振动频率}
%% answer: D
%% memo:
\memoanswer{
质点的振动速度与波的传播速度毫无联系，故 A 错误；如果产生的是横波则质点的振动方向与波的传播的方向垂直，如果产生的是纵波则质点的振动方向与波的传播的方向共线，故 B 错误；质点在一个周期内走过的路程应等于振幅的四倍，与波长无关，故 C 错误；质点的振动频率与波源的振动频率相等，故 D 正确。
}

\item
%% number: 17
%% paperName: 2011年北京理综第 17 题 6 分
%% typeId: 1
%% type: 单选题
%% chapter: 电路
%% point: 动态分析
%% method: 无
%% score: 6
%% degree: 450
%% duplicateId: 0
%% body:
如图所示电路，电源内阻不可忽略。开关 S 闭合后，在变阻器 $R_{0}$ 的滑动端向下滑动的过程中，\xzanswer{A}

\onepicture[w=6cm]{figs/17.png}

\fourchoices[answer=A]
{电压表与电流表的示数都减小}
{电压表与电流表的示数都增大}
{电压表的示数增大，电流表的示数减小}
{电压表的示数减小，电流表的示数增大}
%% answer: A
%% memo:
\memoanswer{
滑动变阻器滑片向下滑动，则连入电路的阻值变小，并联电路的阻值变小，电路总电阻变小，则干路电流变大，电源内阻分得的电压变大，故路端电压变小，即电压表示数变小。干路电流变大，$R_{1}$ 两端电压变大，则 $R_{2}$ 两端电压变小，电流表示数减小，故 A 正确，其余选项错误。
}

\item
%% number: 18
%% paperName: 2011年北京理综第 18 题 6 分
%% typeId: 1
%% type: 单选题
%% chapter: 运动和力
%% point: 牛顿定律中的图象
%% method: 无
%% score: 6
%% degree: 450
%% duplicateId: 0
%% body:
“蹦极”就是跳跃者把一端固定的长弹性绳绑在踝关节等处，从几十米高处跳下的一种极限运动。某人做蹦极运动，所受绳子拉力 $F$ 的大小随时间 $t$ 变化的情况如图所示。将蹦极过程近似为在竖直方向的运动，重力加速度为 $g$。据图可知，此人在蹦极过程中最大加速度约为\xzanswer{B}

\onepicture[w=7cm]{figs/18.png}

\fourchoices[answer=B]
{$g$}
{$2g$}
{$3g$}
{$4g$}
%% answer: B
%% memo:
\memoanswer{
由图可知，最后绳子的拉力恒为 $0.6F_{0}$，即人的重力为 $0.6F_{0}$，而绳子对人的最大拉力为 $1.8F_{0}$，由牛顿第二定律得 $1.8F_{0}-0.6F_{0}=ma_{\max}$，故 $a_{\max}=2g$，故 B 正确。
}

\item
%% number: 19
%% paperName: 2011年北京理综第 19 题 6 分
%% typeId: 1
%% type: 单选题
%% chapter: 电磁感应
%% point: 自感与互感，涡流
%% method: 无
%% score: 6
%% degree: 450
%% duplicateId: 0
%% body:
某同学为了验证断电自感现象，自己找来带铁心的线圈 $L$、小灯泡 A、开关 S 和电池组 $E$，用导线将它们连接成如图所示的电路。检查电路后，闭合开关 S，小灯泡发光；再断开开关 S，小灯泡仅有不显著的延时熄灭现象。虽经多次重复，仍未见老师演示时出现的小灯泡闪亮现象，他冥思苦想找不出原因。你认为最有可能造成小灯泡未闪亮的原因是\xzanswer{C}

\onepicture[w=6cm]{figs/19.png}

\fourchoices[answer=C]
{电源的内阻较大}
{小灯泡电阻偏大}
{线圈电阻偏大}
{线圈的自感系数较大}
%% answer: C
%% memo:
\memoanswer{
小灯泡没有出现闪亮现象是因为断电后电路中的小灯泡两端电压太小。因断电后电路与电源脱离关系，线圈与灯泡组成闭合回路，故电源内阻大小对自感无影响，故 A 错误；若小灯泡电阻偏大，则分得的电压就大，这有助于出现闪亮现象，故 B 错误；若线圈电阻偏大，在自感电动势一定的情况下，线圈内阻上的电压偏大，相应灯泡两端的电压就偏小，这不利于出现闪亮现象，故 C 正确；线圈自感系数越大，产生的自感电动势越大，这有利于闪亮现象的出现，故 D 错误。
}

\item
%% number: 20
%% paperName: 2011年北京理综第 20 题 6 分
%% typeId: 1
%% type: 单选题
%% chapter: 运动和力
%% point: 力学单位制
%% method: 无
%% score: 6
%% degree: 450
%% duplicateId: 0
%% body:
物理关系式不仅反映了物理量之间的关系，也确定了单位间的关系。如关系式 $U=IR$ 既反映了电压、电流和电阻之间的关系，也确定了 V（伏）与 A（安）和 $\Omega$（欧）的乘积等效。现有物理量单位：m（米）、s（秒）、N（牛）、J（焦）、W（瓦）、C（库）、F（法）、A（安）、$\Omega$（欧）和 T（特），由它们组合成的单位都与电压单位 V（伏）等效的是\xzanswer{B}

\fourchoices[answer=B]
{$\UJC$ 和 $\UNC$}
{$\UCF$ 和 $\UTmqs$}
{$\UWA$ 和 $\UCTms$}
{$\UW^{\frac{1}{2}}\cdot\UO^{\frac{1}{2}}$ 和 $\UTAm$}
%% answer: B
%% memo:
\memoanswer{
由电容的定义式 $C=\frac{Q}{U}$，有 $U=\frac{Q}{C}$，在国际单位制中，电量的单位是 $\UC$，电容的单位是 $\UF$，故电压单位 $\UV$ 与 $\UCF$ 等效；由法拉第电磁感应定律知，长为 $L$ 的导体棒在磁场中以速度 $v$ 做垂直切割磁感线运动时，产生的感应电动势为 $E=BLv$，而在国际单位制中，磁感应强度的单位是 $\UT$，长度的单位是 $\Um$，速度的单位是 $\Ums$，故电压单位 $\UV$ 与 $\UTmqs$ 等效，故 B 正确。
}

\item
%% number: 21
%% paperName: 2011年北京理综第 21 题 8 分
%% typeId: 4
%% type: 实验题
%% chapter: 电路
%% point: 多用表的使用
%% method: 无
%% score: 8
%% degree: 650
%% duplicateId: 0
%% body:
用如图所示的多用电表测量电阻，要用到选择开关 K 和两个部件 S、T。请根据下列步骤完成电阻测量：

\begin{steps}
	\item
	旋动部件\tkanswer[1.3cm]{$S$}，使指针对准电流的“0”刻线。
	\item
	将 K 旋转到电阻挡“$\times$100”的位置。
	\item
	将插入“＋”、“－”插孔的表笔短接，旋动部件\tkanswer[1.3cm]{$T$}，使指针对准电阻的\tkanswer[1.8cm]{0 刻线}（填“0 刻线”或“$\infty$ 刻线”）。
	\item
	将两表笔分别与待测电阻相接，发现指针偏转角度过小。为了得到比较准确的测量结果，请从下列选项中挑出合理的步骤，并按\tkanswer[1.4cm]{ADC}的顺序进行操作，再完成读数测量。
\end{steps}

（A）将 K 旋转到电阻挡“$\times$1k”的位置

（B）将 K 旋转到电阻挡“$\times$10”的位置

（C）将两表笔的金属部分分别与被测电阻的两根引线相接

（D）将两表笔短接，旋动合适部件，对电表进行校准

\onepicture[h=4.5cm, align=right]{figs/21.png}

%% answer: 
\jdanswer{
$S$；$T$，0 刻线；ADC
}
%% memo:
\memoanswer{
（1）①$S$ 为机械调零旋钮，应调整它使指针对准电流的“0”刻线；

③将 K 拨到适当的挡位后，要将两表笔短接后旋动 $T$ 进行调零，使指针对准电阻的“0”刻线；

④指针偏转角度过小，说明所测电阻阻值较大，应选用更高的挡位，要再次将两表笔短接进行调零，然后才可以测电阻，故正确选项和正确顺序为 ADC。
}

\item
%% number: 21
%% paperName: 2011年北京理综第 21 题 10 分
%% typeId: 4
%% type: 实验题
%% chapter: 运动和力
%% point: 动量守恒定律
%% method: 无
%% score: 10
%% degree: 450
%% duplicateId: 0
%% body:
如图所示，用“碰撞实验器”可以验证动量守恒定律，即研究两个小球在轨道水平部分碰撞前后的动量关系。

\begin{steps}
	\item
	实验中，直接测定小球碰撞前后的速度是不容易的。但是，可以通过仅测量\tkanswer[1.2cm]{C}（填选项前的符号），间接地解决这个问题。

	（A）小球开始释放高度 $h$

	（B）小球抛出点距地面的高度 $H$

	（C）小球做平抛运动的射程

	\item
	图 \ref{2011北京理综21a} 中 O 点是小球抛出点在地面上的垂直投影。实验时，先让入射球 $m_{1}$ 多次从斜轨上 $S$ 位置静止释放，找到其平均落地点的位置 P，测量平抛射程 OP。然后，把被碰小球 $m_{2}$ 静置于轨道的水平部分，再将入射球 $m_{1}$ 从斜轨上 S 位置静止释放，与小球 $m_{2}$ 相碰，并多次重复。接下来要完成的必要步骤是\tkanswer[1.2cm]{ADE}（填选项前的符号）。

	（A）用天平测量两个小球的质量 $m_{1}$、$m_{2}$

	（B）测量小球 $m_{1}$ 开始释放高度 $h$

	（C）测量抛出点距地面的高度 $H$

	（D）分别找到 $m_{1}$、$m_{2}$ 相碰后平均落地点的位置 M、N

	（E）测量平抛射程 OM，ON

	\item
	若两球相碰前后的动量守恒，其表达式可表示为\tkanswer[3.4cm]{$m_{1}\cdot OM+m_{2}\cdot ON=m_{1}\cdot OP$}（用②中测量的量表示）；

	若碰撞是弹性碰撞，那么还应满足的表达式为\tkanswer[5.8cm]{$m_{1}\cdot OM^{2}+m_{2}\cdot ON^{2}=m_{1}\cdot OP^{2}$}（用②中测量的量表示）。

	\item
	经测定，$m_{1}=45.0\Ug$，$m_{2}=7.5\Ug$，小球落地点的平均位置距 O 点的距离如图 \ref{2011北京理综21b} 所示。碰撞前、后 $m_{1}$ 的动量分别为 $p_{1}$ 与 $p_{1}^{\prime}$，则 $p_{1}:p_{1}^{\prime}=$\tkanswer[1cm]{14}$:11$；若碰撞结束时 $m_{2}$ 的动量为 $p_{2}^{\prime}$，则 $p_{1}^{\prime}:p_{2}^{\prime}=11:$\tkanswer[1cm]{2.9}。

	实验结果表明，碰撞前、后总动量的比值 $\frac{p_{1}}{p_{1}^{\prime}+p_{2}^{\prime}}$ 为\tkanswer[1.6cm]{$1\sim1.01$}。

	\item
	有同学认为，在上述实验中仅更换两个小球的材质，其它条件不变，可以使被碰小球做平抛运动的射程增大。请你用④中已知的数据，分析和计算出被碰小球 $m_{2}$ 平抛运动射程 ON 的最大值为\tkanswer[1.4cm]{76.8}$\Ucm$。
\end{steps}

\twopicture[labelA=2011北京理综21a, labelB=2011北京理综21b, numA=1, numB=2, align=right, wA=5.5cm, wB=6cm]
{\includesvg{figs/21a}}
{figs/21b.png}

%% answer: 
\jdanswer{
①C；②ADE 或 DEA 或 DAE；③$m_{1}\cdot OM+m_{2}\cdot ON=m_{1}\cdot OP$，$m_{1}\cdot OM^{2}+m_{2}\cdot ON^{2}=m_{1}\cdot OP^{2}$；④14，2.9，$1\sim1.01$；⑤76.8
}
%% memo:
\memoanswer{
（2）①因为小球离开桌面后做平抛运动，水平方向上做匀速直线运动，而竖直方向下落高度相同，所用时间相同，水平方向的位移之比等于速度之比，故 C 正确。

②由以上分析再结合动量守恒定律可知 ADE 正确。

③由动量守恒定律得 $m_{1}v_{1}+m_{2}v_{2}=m_{1}v$，即 $m_{1}\cdot OM+m_{2}\cdot ON=m_{1}\cdot OP$；若是弹性碰撞，碰撞前后两个小球组成的系统动能相等，则 $\frac{1}{2}m_{1}v_{1}^{2}+\frac{1}{2}m_{2}v_{2}^{2}=\frac{1}{2}m_{1}v^{2}$，即 $m_{1}\cdot OM^{2}+m_{2}\cdot ON^{2}=m_{1}\cdot OP^{2}$。

④由 $\frac{p_{1}}{p_{1}^{\prime}}=\frac{OP}{OM}=\frac{44.80}{35.20}=\frac{14}{11}$，$\frac{p_{1}^{\prime}}{p_{2}^{\prime}}=\frac{m_{1}\cdot OM}{m_{2}\cdot ON}=\frac{45.0\times35.20}{7.5\times55.68}=\frac{11}{2.9}$，$\frac{p_{1}}{p_{1}^{\prime}+p_{2}^{\prime}}=\frac{m_{1}\cdot OP}{m_{1}\cdot OM+m_{2}\cdot ON}=\frac{45.0\times44.80}{45.0\times35.20+7.5\times55.68}\approx1.007$。

⑤只改变小球的材质而其他条件不变，要使 $m_{2}$ 水平射程最大，则两球应发生完全弹性碰撞，即碰撞前后动量守恒且动能无损失，有 $m_{1}\cdot OM+m_{2}\cdot ON=m_{1}\cdot OP$，$m_{1}\cdot OM^{2}+m_{2}\cdot ON^{2}=m_{1}\cdot OP^{2}$，其中 $m_{1}=45.0\Ug$，$m_{2}=7.5\Ug$，$OP=44.80\Ucm$，解得 $ON=76.8\Ucm$。
}

\item
%% number: 22
%% paperName: 2011年北京理综第 22 题 16 分
%% typeId: 6
%% type: 计算题
%% chapter: 曲线运动
%% point: 竖直平面圆周运动
%% method: 无
%% score: 16
%% degree: 650
%% duplicateId: 0
%% body:
如图所示，长度为 $l$ 的轻绳上端固定在 $O$ 点，下端系一质量为 $m$ 的小球（小球的大小可以忽略）。

\begin{enumerate}
	\item
	在水平拉力 $F$ 的作用下，轻绳与竖直方向的夹角为 $\alpha$，小球保持静止。画出此时小球的受力图，并求力 $F$ 的大小；
	\item
	由图示位置无初速释放小球，求当小球通过最低点时的速度大小及轻绳对小球的拉力。不计空气阻力。
\end{enumerate}

\onepicture[align=right, h=4.5cm]{figs/22.png}

%% answer: 
\jdanswer{
\begin{enumerate}
	\item
	$F=mg\tan\alpha$
	\item
	$v=\sqrt{2gl(1-\cos\alpha)}$，$F_{\mathrm{T}}^{\prime}=mg(3-2\cos\alpha)$，方向竖直向上
\end{enumerate}
}
%% memo:
\memoanswer{
（1）受力如图所示。

\onepicture[w=3.2cm]{figs/22_an.jpg}

由平衡条件得 $F_{\mathrm{T}}\cos\alpha=mg$，$F_{\mathrm{T}}\sin\alpha=F$，联立解得 $F=mg\tan\alpha$。

（2）运动中只有重力做功，系统机械能守恒得 $mgl(1-\cos\alpha)=\frac{1}{2}mv^{2}$，则通过最低点时，小球的速度大小 $v=\sqrt{2gl(1-\cos\alpha)}$，由牛顿第二定律得 $F_{\mathrm{T}}^{\prime}-mg=m\frac{v^{2}}{l}$，解得轻绳对小球的拉力 $F_{\mathrm{T}}^{\prime}=mg(3-2\cos\alpha)$，方向竖直向上。
}

\item
%% number: 23
%% paperName: 2011年北京理综第 23 题 18 分
%% typeId: 6
%% type: 计算题
%% chapter: 磁场
%% point: 带电粒子在组合场中的运动
%% method: 无
%% score: 18
%% degree: 450
%% duplicateId: 0
%% body:
利用电场和磁场，可以将比荷不同的离子分开，这种方法在化学分析和原子核技术等领域有重要的应用。

如图所示的矩形区域 ACDG（AC 边足够长）中存在垂直于纸面的匀强磁场，A 处有一狭缝。离子源产生的离子，经静电场加速后穿过狭缝沿垂直于 GA 边且垂直于磁场的方向射入磁场，运动到 GA 边，被相应的收集器收集。整个装置内部为真空。

已知被加速的两种正离子的质量分别是 $m_{1}$ 和 $m_{2}$（$m_{1}>m_{2}$），电荷量均为 $q$。加速电场的电势差为 $U$，离子进入电场时的初速度可以忽略。不计重力，也不考虑离子间的相互作用。

\begin{enumerate}
	\item
	求质量为 $m_{1}$ 的离子进入磁场时的速率 $v_{1}$；
	\item
	当磁感应强度的大小为 $B$ 时，求两种离子在 GA 边落点的间距 $s$；
	\item
	在前面的讨论中忽略了狭缝宽度的影响，实际装置中狭缝具有一定宽度。若狭缝过宽，可能使两束离子在 GA 边上的落点区域交叠，导致两种离子无法完全分离。

	设磁感应强度大小可调，GA 边长为定值 $L$，狭缝宽度为 $d$，狭缝右边缘在 A 处。离子可以从狭缝各处射入磁场，入射方向仍垂直于 GA 边且垂直于磁场。为保证上述两种离子能落在 GA 边上并被完全分离，求狭缝的最大宽度。
\end{enumerate}

\onepicture[align=right, w=6.5cm]{figs/23.png}

%% answer: 
\jdanswer{
\begin{enumerate}
	\item
	$v_{1}=\sqrt{\frac{2qU}{m_{1}}}$
	\item
	$s=\sqrt{\frac{8U}{qB^{2}}}(\sqrt{m_{1}}-\sqrt{m_{2}})$
	\item
	$d_{m}=\frac{\sqrt{m_{1}}-\sqrt{m_{2}}}{2\sqrt{m_{1}}-\sqrt{m_{2}}}L$
\end{enumerate}
}
%% memo:
\memoanswer{
（1）加速电场对离子 $m_{1}$ 做的功为 $W=qU$，由动能定理得 $\frac{1}{2}m_{1}v_{1}^{2}=qU$，解得 $v_{1}=\sqrt{\frac{2qU}{m_{1}}}$①。

（2）由牛顿第二定律和洛伦兹力提供向心力得 $qvB=\frac{mv^{2}}{R}$，解得 $R=\frac{mv}{qB}$，利用①式得离子在磁场中的轨道半径分别为 $R_{1}=\sqrt{\frac{2m_{1}U}{qB^{2}}}$，$R_{2}=\sqrt{\frac{2m_{2}U}{qB^{2}}}$②。

两种离子在 GA 上落点的间距为 $s=2R_{1}-2R_{2}=\sqrt{\frac{8U}{qB^{2}}}(\sqrt{m_{1}}-\sqrt{m_{2}})$③。

（3）质量为 $m_{1}$ 的离子，在 GA 边上的落点都在其入射点左侧 $2R_{1}$ 处，由于狭缝的宽度为 $d$，因此落点区域的宽度也是 $d$。同理，质量为 $m_{2}$ 的离子在 GA 边上落点区域的宽度也是 $d$。

为保证两种离子能完全分离，两个区域应无交叠，条件为 $2(R_{1}-R_{2})>d$④。

利用②式，代入④式得 $2R_{1}\left(1-\sqrt{\frac{m_{2}}{m_{1}}}\right)>d$，$R_{1}$ 的最大值满足 $2R_{1m}=L-d$，得 $(L-d)\left(1-\sqrt{\frac{m_{2}}{m_{1}}}\right)>d$，求得最大值 $d_{m}=\frac{\sqrt{m_{1}}-\sqrt{m_{2}}}{2\sqrt{m_{1}}-\sqrt{m_{2}}}\cdot L$。
}

\item
%% number: 24
%% paperName: 2011年北京理综第 24 题 20 分
%% typeId: 6
%% type: 计算题
%% chapter: 电场
%% point: 带电粒子的往复运动
%% method: 无
%% score: 20
%% degree: 450
%% duplicateId: 0
%% body:
静电场方向平行于 $x$ 轴，其电势 $\varphi$ 随 $x$ 的分布可简化为如图所示的折线，图中 $\varphi_{0}$ 和 $d$ 为已知量。一个带负电的粒子在电场中以 $x=0$ 为中心，沿 $x$ 轴方向做周期性运动。已知该粒子质量为 $m$、电量为 $-q$，其动能与电势能之和为 $-A$（$0<A<q\varphi_{0}$）。忽略重力。求：

\begin{enumerate}
	\item
	粒子所受电场力的大小；
	\item
	粒子的运动区间；
	\item
	粒子的运动周期。
\end{enumerate}

\onepicture[align=right, w=6cm]{figs/24.png}

%% answer: 
\jdanswer{
\begin{enumerate}
	\item
	$F=qE=\frac{q\varphi_{0}}{d}$
	\item
	$-d\left(1-\frac{A}{q\varphi_{0}}\right)\leq x\leq d\left(1-\frac{A}{q\varphi_{0}}\right)$
	\item
	$T=4t=\frac{4d}{q\varphi_{0}}\sqrt{2m(q\varphi_{0}-A)}$
\end{enumerate}
}
%% memo:
\memoanswer{
（1）由图可知，$O$ 与 $d$（或 $-d$）两点间的电势差为 $\varphi_{0}$，电场强度的大小 $E=\frac{\varphi_{0}}{d}$，电场力的大小为 $F=qE=\frac{q\varphi_{0}}{d}$。

（2）设粒子在 $[-x_{0},x_{0}]$ 区间内运动，速率为 $v$，由题意得 $\frac{1}{2}mv^{2}-q\varphi=-A$①，由图可知 $\varphi=\varphi_{0}\left(1-\frac{|x|}{d}\right)$②，由①②得 $\frac{1}{2}mv^{2}=q\varphi_{0}\left(1-\frac{|x|}{d}\right)-A$③，因动能非负，有 $q\varphi_{0}\left(1-\frac{|x|}{d}\right)-A\geq0$，得 $|x|\leq d\left(1-\frac{A}{q\varphi_{0}}\right)$，即 $x_{0}=d\left(1-\frac{A}{q\varphi_{0}}\right)$④，粒子的运动区间为 $-d\left(1-\frac{A}{q\varphi_{0}}\right)\leq x\leq d\left(1-\frac{A}{q\varphi_{0}}\right)$。

（3）考虑粒子从 $-x_{0}$ 处开始运动的四分之一周期，由牛顿第二定律得，粒子的加速度 $a=\frac{F}{m}=\frac{qE}{m}=\frac{q\varphi_{0}}{md}$⑤，由匀加速直线运动得 $t=\sqrt{\frac{2x_{0}}{a}}$，将④⑤代入，得 $t=\sqrt{\frac{2md^{2}}{q\varphi_{0}}\left(1-\frac{A}{q\varphi_{0}}\right)}$，粒子的运动周期 $T=4t=\frac{4d}{q\varphi_{0}}\sqrt{2m(q\varphi_{0}-A)}$⑥。
}

\end{enumerate}

\end{document}
